% dsh-latex showcase — compile with:
%   latex_compile { "source_path": "examples/showcase.tex" }
% or from a shell:
%   node <dsh-latex>/scripts/warm-cache.mjs --profile=chinese
%   <runtimeDir>/tectonic -X compile examples/showcase.tex --outdir examples
\documentclass[11pt]{ctexart}

\usepackage{amsmath,amssymb}
\usepackage{booktabs}
\usepackage{geometry}
\usepackage{hyperref}

\title{dsh-latex 示例文档}
\author{DeepSeek Harness}
\date{\today}

\begin{document}
\maketitle

\section{这是什么}
这是 \texttt{dsh-latex} 的示例文档：由固定版本的 Tectonic 引擎编译，
不需要系统安装 TeX，编译时只按需拉取文档真正用到的宏包。

\section{数学}
行内公式 $E = mc^2$，以及行间公式：
\begin{equation}
  \int_{-\infty}^{\infty} e^{-x^{2}}\,\mathrm{d}x = \sqrt{\pi}
  \qquad\text{和}\qquad
  \sum_{n=1}^{\infty}\frac{1}{n^{2}} = \frac{\pi^{2}}{6}.
\end{equation}

矩阵与对齐环境同样可用：
\begin{align}
  \nabla \cdot \mathbf{E} &= \frac{\rho}{\varepsilon_0}, &
  \nabla \times \mathbf{B} &= \mu_0 \mathbf{J}
  + \mu_0\varepsilon_0 \frac{\partial \mathbf{E}}{\partial t}.
\end{align}

\section{表格}
\begin{table}[ht]
  \centering
  \begin{tabular}{lrr}
    \toprule
    工具 & 输入 & 输出 \\
    \midrule
    \texttt{latex\_compile} & \texttt{.tex} 文件或内联源码 & PDF \\
    \texttt{latex\_math}    & 单个公式                  & 裁剪后的 PDF \\
    \texttt{latex\_health}  & 无                        & 引擎与缓存状态 \\
    \bottomrule
  \end{tabular}
  \caption{三个工具。}
\end{table}

\section{中文}
中文排版由 \texttt{ctex} 提供（对应 \texttt{warm-cache.mjs --profile=chinese}）。
预热的宏包集合之外的内容，首次编译时会联网补齐，之后即缓存复用。

\end{document}
